% CORRECTED VERSION
% Changes:
% - Fixed the incorrect bound on the noise‑variance sum (Step 10) and the resulting
%   convergence rate (Step 11).
% - Strengthened the step‑size condition to 0<α≤1/(2L) and derived a bound for the
%   *averaged* iterate, which yields the provable O((log K)/K) rate under the
%   original noise schedule σ_k² ≤ c/(k+1).
% - Added Lemma \ref{lem:avg-convex} and an explicit averaging argument.
% - Updated Theorem \ref{thm:main} accordingly while preserving the rest of the
%   document structure.

\documentclass[11pt]{article}
\usepackage{amsmath,amssymb,amsthm}
\usepackage{bm}
\usepackage{algorithm}
\usepackage{algpseudocode}
\usepackage{hyperref}
\usepackage{geometry}
\geometry{margin=1in}
\newtheorem{theorem}{Theorem}
\newtheorem{lemma}{Lemma}
\newtheorem{assumption}{Assumption}
\newcommand{\EE}{\mathbb{E}}
\newcommand{\RR}{\mathbb{R}}
\newcommand{\norm}[1]{\left\lVert #1\right\rVert}
\newcommand{\inner}[1]{\left\langle #1\right\rangle}
\begin{document}

\section*{Algorithm Name}
\textbf{Stochastic Gradient Langevin Dynamics with Decreasing Noise (SGLD‑Dec)}  

The method performs a gradient descent step on a convex smooth objective and adds isotropic Gaussian noise whose variance decays with the iteration index.

\section*{Mathematical Formulation}
Let $f:\RR^{d}\rightarrow\RR$ be the objective function.  
Given an initial point $x_{0}\in\RR^{d}$, step size $\alpha>0$ and a non‑increasing sequence of noise variances $\{\sigma_{k}^{2}\}_{k\ge 0}$, the iterates are defined by  
\begin{equation}
\boxed{%
x_{k+1}=x_{k}-\alpha\,\nabla f(x_{k})+\xi_{k},\qquad 
\xi_{k}\sim\mathcal N\!\big(0,\sigma_{k}^{2}I_{d}\big),\;k=0,1,\dots
}
\label{eq:update}
\end{equation}
The random vectors $\xi_{k}$ are assumed independent of each other and of the history $\{x_{t}\}_{t\le k}$ conditional on $x_{k}$.

\section*{Pseudocode}
\begin{algorithm}[H]
\caption{SGLD‑Dec}
\begin{algorithmic}[1]
\State \textbf{Input:} $x_{0}\in\RR^{d}$, step size $\alpha>0$, variance schedule $\{\sigma_{k}^{2}\}$
\For{$k=0$ \textbf{to} $K-1$}
    \State Compute (full) gradient $g_{k}\gets\nabla f(x_{k})$
    \State Sample noise $\xi_{k}\sim\mathcal N(0,\sigma_{k}^{2}I_{d})$
    \State Update $x_{k+1}\gets x_{k}-\alpha g_{k}+\xi_{k}$ \label{alg:update}
\EndFor
\State \textbf{Output:} $\displaystyle \bar{x}_{K}\;:=\;\frac{1}{K}\sum_{k=0}^{K-1}x_{k}$ \quad (averaged iterate)
\end{algorithmic}
\end{algorithm}

\section*{Dry Run Example}
Consider the one‑dimensional quadratic $f(w)=(w-3)^{2}$.  
\[
\nabla f(w)=2(w-3).
\]
Choose $\alpha=0.1$, variance schedule $\sigma_{k}^{2}=0$ (deterministic walk for illustration), and $x_{0}=0$.

\begin{center}
\begin{tabular}{c|c|c|c}
Iteration $k$ & $x_{k}$ & $\nabla f(x_{k})$ & $x_{k+1}=x_{k}-\alpha\nabla f(x_{k})$ \\ \hline
0 & $0$ & $2(0-3)=-6$ & $0-0.1(-6)=0.6$ \\
1 & $0.6$ & $2(0.6-3)=-4.8$ & $0.6-0.1(-4.8)=1.08$ \\
2 & $1.08$ & $2(1.08-3)=-3.84$ & $1.08-0.1(-3.84)=1.464$ \\
\end{tabular}
\end{center}
Corresponding objective values:
\[
f(0)=9,\; f(0.6)=5.76,\; f(1.08)=3.6864,\; f(1.464)=2.340\ldots
\]
The iterates move toward the minimiser $w^{*}=3$.

\newpage
\section*{Assumptions}
\begin{assumption}[Convexity and Smoothness]\label{ass:convex-smooth}
The function $f$ is convex and $L$‑smooth, i.e.
\[
f(y)\ge f(x)+\inner{\nabla f(x)}{y-x},\qquad 
\|\nabla f(x)-\nabla f(y)\|\le L\|x-y\|\quad\forall x,y\in\RR^{d}.
\]
\end{assumption}
\begin{assumption}[Existence of a Minimiser]\label{ass:opt}
There exists $x^{\star}\in\RR^{d}$ such that $f(x^{\star})=\min_{x}f(x)$ and $\nabla f(x^{\star})=0$.
\end{assumption}
\begin{assumption}[Noise Schedule]\label{ass:noise}
The noise vectors $\{\xi_{k}\}$ satisfy  
\[
\EE[\xi_{k}\mid x_{k}]=0,\qquad
\EE\big[\|\xi_{k}\|^{2}\mid x_{k}\big]=d\,\sigma_{k}^{2},
\]
with a deterministic, non‑increasing sequence $\sigma_{k}^{2}$ such that
\[
\sigma_{k}^{2}\le \frac{c}{k+1}\quad\text{for some }c>0,\;k\ge0.
\tag{A3}
\]
\end{assumption}
\begin{assumption}[Bounded Initial Distance]\label{ass:init}
Denote $D_{0}^{2}:=\|x_{0}-x^{\star}\|^{2}<\infty$.
\end{assumption}

\section*{Main Convergence Theorem}
\begin{theorem}[Expected Suboptimality of the Averaged Iterate]\label{thm:main}
Under Assumptions \ref{ass:convex-smooth}–\ref{ass:init} and for a constant step size
\[
0<\alpha\le\frac{1}{2L},
\]
the averaged iterate $\displaystyle \bar{x}_{K}:=\frac{1}{K}\sum_{k=0}^{K-1}x_{k}$
generated by \eqref{eq:update} satisfies for every $K\ge1$
\[
\boxed{%
\EE\!\big[f(\bar{x}_{K})-f(x^{\star})\big]\;\le\;
\frac{\norm{x_{0}-x^{\star}}^{2}}{\alpha K}
\;+\;
\frac{d\,c\big(1+\log K\big)}{\alpha K}\,.
}
\tag{T1}
\]
Consequently, $\EE[f(\bar{x}_{K})-f(x^{\star})]=\mathcal O\!\big(\tfrac{\log K}{K}\big)$.
\end{theorem}

\section*{Theoretical Properties}
\begin{itemize}
\item \textbf{Stability.} The condition $\alpha\le 1/(2L)$ guarantees that the deterministic part (gradient descent) is contractive enough to dominate the stochastic perturbation.
\item \textbf{Hyperparameter Sensitivity.} A larger $\alpha$ reduces the first term in \eqref{T1} but also enlarges the coefficient of the stochastic term; the optimal admissible choice is $\alpha=1/(2L)$.
\item \textbf{Comparison to Baselines.}
  \begin{enumerate}
    \item Pure GD ($\xi_{k}=0$) yields $\EE[f(\bar{x}_{K})-f(x^{\star})]\le \frac{\norm{x_{0}-x^{\star}}^{2}}{\alpha K}$.
    \item Classical SGLD with constant variance $\sigma^{2}>0$ only guarantees a steady‑state error $\mathcal O(\sigma^{2})$; the decreasing schedule reduces this bias to $\mathcal O\big(\tfrac{\log K}{K}\big)$.
  \end{enumerate}
\end{itemize}

\section*{Discussion}
The theorem shows that, even with isotropic Gaussian perturbations that decay only as $1/k$, the algorithm retains a near‑optimal $\mathcal O(1/K)$ rate up to a logarithmic factor when we consider the averaged iterate. The extra $\log K$ term is unavoidable under the harmonic decay of the variance; a faster decay (e.g. $\sigma_{k}^{2}=O(1/k^{2})$) would eliminate it.

\newpage
\section*{Preliminaries (Definitions and Lemmas)}
\begin{lemma}[Smoothness Inequality]\label{lem:smooth}
If $f$ is $L$‑smooth then for any $x,y\in\RR^{d}$
\[
f(y)\le f(x)+\inner{\nabla f(x)}{y-x}+\frac{L}{2}\|y-x\|^{2}.
\tag{L1}
\]
\end{lemma}
\begin{proof}
Standard consequence of the definition of $L$‑smoothness; see e.g. Nesterov (2004). \hfill $\square$
\end{proof}

\begin{lemma}[Gradient Lower Bound for Convex $L$‑smooth functions]\label{lem:grad-bound}
For convex $L$‑smooth $f$ and any $x$,
\[
\frac{1}{2L}\|\nabla f(x)\|^{2}\le f(x)-f(x^{\star}).
\tag{L2}
\]
\end{lemma}
\begin{proof}
Combine convexity with smoothness (see, e.g., Lemma 1.2.3 in \cite{Bubeck2015}). \hfill $\square$
\end{proof}

\begin{lemma}[Averaging for Convex Functions]\label{lem:avg-convex}
If $f$ is convex and $\bar{x}_{K}:=\frac{1}{K}\sum_{k=0}^{K-1}x_{k}$, then
\[
f(\bar{x}_{K})-f(x^{\star})\le\frac{1}{K}\sum_{k=0}^{K-1}\bigl(f(x_{k})-f(x^{\star})\bigr).
\tag{L3}
\]
\end{lemma}
\begin{proof}
Convexity implies $f\bigl(\frac{1}{K}\sum x_{k}\bigr)\le\frac{1}{K}\sum f(x_{k})$. Subtract $f(x^{\star})$ from both sides. \hfill $\square$
\end{proof}

\section*{Main Proof (Step‑by‑Step Derivation)}
All expectations are taken with respect to the randomness up to iteration $k$, i.e. conditional on $x_{k}$ where appropriate.

\bigskip
\noindent\textbf{Step 1. Distance recursion.}  
From the update \eqref{eq:update},
\[
\begin{aligned}
\|x_{k+1}-x^{\star}\|^{2}
&=\|x_{k}-\alpha\nabla f(x_{k})+\xi_{k}-x^{\star}\|^{2}\\
&=\|x_{k}-x^{\star}\|^{2}
-2\alpha\inner{\nabla f(x_{k})}{x_{k}-x^{\star}}
+ \alpha^{2}\|\nabla f(x_{k})\|^{2}\\
&\qquad
+2\inner{\xi_{k}}{x_{k}-x^{\star}-\alpha\nabla f(x_{k})}
+\|\xi_{k}\|^{2}. \tag{S1}
\end{aligned}
\]

\noindent\textbf{Step 2. Conditional expectation.}  
Using $\EE[\xi_{k}\mid x_{k}]=0$ and $\EE[\|\xi_{k}\|^{2}\mid x_{k}]=d\sigma_{k}^{2}$,
\[
\EE\!\big[\|x_{k+1}-x^{\star}\|^{2}\mid x_{k}\big]
=
\|x_{k}-x^{\star}\|^{2}
-2\alpha\inner{\nabla f(x_{k})}{x_{k}-x^{\star}}
+ \alpha^{2}\|\nabla f(x_{k})\|^{2}
+d\sigma_{k}^{2}. \tag{S2}
\]

\noindent\textbf{Step 3. Apply convexity.}  
Convexity gives $\inner{\nabla f(x_{k})}{x_{k}-x^{\star}}\ge f(x_{k})-f(x^{\star})$. Substituting into \eqref{S2} yields
\[
\EE\!\big[\|x_{k+1}-x^{\star}\|^{2}\mid x_{k}\big]
\le
\|x_{k}-x^{\star}\|^{2}
-2\alpha\bigl(f(x_{k})-f(x^{\star})\bigr)
+ \alpha^{2}\|\nabla f(x_{k})\|^{2}
+d\sigma_{k}^{2}. \tag{S3}
\]

\noindent\textbf{Step 4. Bound the gradient term using Lemma \ref{lem:grad-bound}.}  
From \eqref{L2},
\[
\|\nabla f(x_{k})\|^{2}\le 2L\bigl(f(x_{k})-f(x^{\star})\bigr).
\]
Hence, with $\alpha\le\frac{1}{2L}$,
\[
\alpha^{2}\|\nabla f(x_{k})\|^{2}
\le 2\alpha^{2}L\bigl(f(x_{k})-f(x^{\star})\bigr)
\le \alpha\bigl(f(x_{k})-f(x^{\star})\bigr). \tag{S4}
\]
% CORRECTED: the previous proof cancelled the gradient term completely, which prevented a useful summation. We keep a residual $\alpha$‑scaled term.

\noindent\textbf{Step 5. Insert \eqref{S4} into \eqref{S3}.}  
\[
\EE\!\big[\|x_{k+1}-x^{\star}\|^{2}\mid x_{k}\big]
\le
\|x_{k}-x^{\star}\|^{2}
-\alpha\bigl(f(x_{k})-f(x^{\star})\bigr)
+d\sigma_{k}^{2}. \tag{S5}
\]

\noindent\textbf{Step 6. Rearrange to isolate the suboptimality term.}  
\[
\alpha\bigl(f(x_{k})-f(x^{\star})\bigr)
\le
\|x_{k}-x^{\star}\|^{2}
-\EE\!\big[\|x_{k+1}-x^{\star}\|^{2}\mid x_{k}\big]
+d\sigma_{k}^{2}. \tag{S6}
\]

\noindent\textbf{Step 7. Take total expectation.}  
\[
\alpha\,\EE\!\big[f(x_{k})-f(x^{\star})\big]
\le
\EE\!\big[\|x_{k}-x^{\star}\|^{2}\big]
-\EE\!\big[\|x_{k+1}-x^{\star}\|^{2}\big]
+d\sigma_{k}^{2}. \tag{S7}
\]

\noindent\textbf{Step 8. Sum over $k=0,\dots,K-1$.}  
Telescoping the first two terms gives
\[
\alpha\sum_{k=0}^{K-1}\EE\!\big[f(x_{k})-f(x^{\star})\big]
\le
\|x_{0}-x^{\star}\|^{2}
+d\sum_{k=0}^{K-1}\sigma_{k}^{2}. \tag{S8}
\]

\noindent\textbf{Step 9. Bound the noise‑variance sum using Assumption \ref{ass:noise}.}  
Since $\sigma_{k}^{2}\le c/(k+1)$,
\[
\sum_{k=0}^{K-1}\sigma_{k}^{2}
\le
c\sum_{k=1}^{K}\frac{1}{k}
\le
c\bigl(1+\log K\bigr). \tag{S9}
\]
% CORRECTED: previously the proof incorrectly claimed the sum is bounded by a constant.

\noindent\textbf{Step 10. Divide by $K$ and use Lemma \ref{lem:avg-convex}.}  
Define the averaged iterate $\displaystyle \bar{x}_{K}:=\frac{1}{K}\sum_{k=0}^{K-1}x_{k}$.  
By Lemma \ref{lem:avg-convex},
\[
\EE\!\big[f(\bar{x}_{K})-f(x^{\star})\big]
\le
\frac{1}{K}\sum_{k=0}^{K-1}\EE\!\big[f(x_{k})-f(x^{\star})\big].
\]
Combining this with \eqref{S8}–\eqref{S9} we obtain
\[
\EE\!\big[f(\bar{x}_{K})-f(x^{\star})\big]
\le
\frac{\|x_{0}-x^{\star}\|^{2}}{\alpha K}
\;+\;
\frac{d\,c\big(1+\log K\big)}{\alpha K}. \tag{S10}
\]

\noindent\textbf{Step 11. Conclusion.}  
Inequality \eqref{S10} is precisely the bound stated in Theorem \ref{thm:main}. \hfill $\square$

\section*{Rate Derivation (With Constants)}
From \eqref{S10} we read the explicit constants:
\[
C_{1}:=\frac{\norm{x_{0}-x^{\star}}^{2}}{\alpha},\qquad
C_{2}:=\frac{d\,c\big(1+\log K\big)}{\alpha},
\]
so that
\[
\EE\!\big[f(\bar{x}_{K})-f(x^{\star})\big]\le \frac{C_{1}+C_{2}}{K}
= \mathcal O\!\Big(\frac{\log K}{K}\Big).
\]

\section*{Special Cases}
\begin{enumerate}
\item \textbf{Zero Noise.} $\sigma_{k}^{2}=0$ for all $k$ reduces \eqref{S10} to the classic GD bound $\EE[f(\bar{x}_{K})-f(x^{\star})]\le \frac{\|x_{0}-x^{\star}\|^{2}}{\alpha K}$.
\item \textbf{Constant Noise.} If $\sigma_{k}^{2}\equiv \sigma^{2}$, the sum in \eqref{S8} grows linearly with $K$, leading to a steady‑state error $\EE[f(\bar{x}_{K})-f(x^{\star})]\le \frac{d\sigma^{2}}{\alpha}$, i.e. no $1/K$ decay.
\item \textbf{Stronger Decay of Noise.} If the schedule satisfies $\sigma_{k}^{2}\le c/(k+1)^{2}$, then $\sum_{k=0}^{K-1}\sigma_{k}^{2}\le c$, and \eqref{S10} simplifies to the pure $O(1/K)$ rate without the logarithmic factor.
\item \textbf{Strongly Convex $f$.} If additionally $f$ is $\mu$‑strongly convex, a similar derivation yields an exponential term plus a $\frac{d c}{\mu K}$ stochastic contribution.
\end{enumerate}

\bigskip
\noindent\textbf{References}
\begin{enumerate}
\item Bubeck, S. (2015). \emph{Convex Optimization: Algorithms and Complexity}. Foundations and Trends in Machine Learning.
\item Welling, M., \& Teh, Y. W. (2011). Bayesian learning via stochastic gradient Langevin dynamics. \emph{Proceedings of the 28th International Conference on Machine Learning (ICML-11)}.
\end{enumerate}

\end{document}

% ===== CORRECTION SUMMARY =====
% 1. Step 10 Issue: The original proof incorrectly bounded $\sum_{k=0}^{K-1}\sigma_k^2$ by a constant. Fixed by applying the harmonic series bound $\sum_{k=1}^{K}1/k \le 1+\log K$ (see S9) and carrying the $\log K$ term through the analysis.
% 2. Step 11 Issue: The division by $K$ that turned a constant into a $1/K$ term was unjustified. Updated the theorem to bound the *averaged* iterate (Lemma L3) and derived the correct $\mathcal O((\log K)/K)$ rate (see S10–S11). 
% 3. Additional adjustments: strengthened the step‑size condition to $0<\alpha\le 1/(2L)$, added Lemma \ref{lem:avg-convex} for averaging, and rewrote the recursion to retain a non‑vanishing gradient term (S4–S5).